% Source: https://artofproblemsolving.com/community/c3h2274477_counting
% Stable identity: aops-topic-2274477-q01
\textbf{(c)} The general rule first. The $k$th number in the list is $a+(k-1)d$. If $b$ is the $n$th number, then $a+(n-1)d=b$, so
\[n=\dfrac{b-a}{d}+1.\]
In words: $\dfrac{b-a}{d}$ counts the \emph{steps} (gaps) from $a$ to $b$, and there is one more number than there are gaps (like fence posts and the spaces between them). Forgetting the ``$+1$'' is the usual mistake.

\textbf{(a)} Here $d=1$: $97-25+1=\boxed{73}$ integers. (Equivalently, subtract $24$ from every number to get the list $1,2,\ldots,73$.)

\textbf{(b)} Here $d=4$: $\dfrac{403-7}{4}+1=99+1=\boxed{100}$ numbers.
